\documentclass[12pt,a4paper]{article}

% ============================================================
% PACKAGES
% ============================================================
\usepackage[left=1in,right=1in,top=1in,bottom=1in]{geometry}
\usepackage{graphicx}
\usepackage{booktabs}
\usepackage{array}
\usepackage{tabularx}
\usepackage{colortbl}
\usepackage{xcolor}
\usepackage{makecell}
\usepackage{fancyhdr}
\usepackage{titlesec}
\usepackage{amsmath}
\usepackage{tcolorbox}

\definecolor{headerblue}{RGB}{213,220,228}

\newcolumntype{Y}{>{\centering\arraybackslash}X}
\newcolumntype{L}{>{\raggedright\arraybackslash}X}

\setlength{\parindent}{0pt}
\setlength{\parskip}{6pt}

\pagestyle{fancy}
\fancyhf{}
\rhead{SE: Machine Learning Lab}
\lhead{Lab 6}
\rfoot{\thepage}

\begin{document}

% ============================================================
% TABLE 1: DEPARTMENT HEADER
% ============================================================
\renewcommand{\arraystretch}{2.0}
\begin{tabularx}{\textwidth}{|Y|}
\hline
\rowcolor{headerblue} \textbf{\large Department of Computer and Software Engineering} \\
\hline
\textbf{\large SE: Machine Learning} \\
\hline
\end{tabularx}
\renewcommand{\arraystretch}{1.0}

\vspace{12pt}

% ============================================================
% TABLE 2: INSTRUCTOR INFO
% ============================================================
\renewcommand{\arraystretch}{2.0}
\begin{tabularx}{\textwidth}{|L|L|}
\hline
\textbf{Course Instructor:} Dr. Ahmed Raza & \textbf{Date:} \\
\hline
\textbf{Day:} & \textbf{Semester:} \\
\hline
\end{tabularx}
\renewcommand{\arraystretch}{1.0}

\vspace{12pt}

% ============================================================
% TITLE
% ============================================================
\begin{center}
{\Large\bfseries Lab 6: K-Means Clustering from Scratch}
\end{center}

\vspace{6pt}

% ============================================================
% TABLE 3: LAB METADATA
% ============================================================
\renewcommand{\arraystretch}{1.8}
\begin{tabularx}{\textwidth}{|Y|c|c|c|c|}
\hline
\textbf{Lab Title} & \textbf{CLO} & \textbf{BT} & \textbf{PLO} & \textbf{Weightage} \\
\hline
\makecell{{\small K-Means Clustering} \\ {\small From Scratch}} & & & & \\
\hline
\end{tabularx}
\renewcommand{\arraystretch}{1.0}

\vspace{6pt}

% ============================================================
% TABLE 4: STUDENT INFO
% ============================================================
\renewcommand{\arraystretch}{2.0}
\begin{tabularx}{\textwidth}{|Y|Y|Y|Y|Y|}
\hline
\textbf{Name} & \textbf{Roll Number} & \textbf{Report /10} & \textbf{Viva /5} & \textbf{Total /15} \\
\hline
 & & & & \\
\hline
\end{tabularx}
\renewcommand{\arraystretch}{1.0}

\vspace{12pt}

\begin{flushright}
Checked on: \_\_\_\_\_\_\_\_\_\_\_\_\_\_\_\_

\vspace{6pt}

Signature: \_\_\_\_\_\_\_\_\_\_\_\_\_\_\_\_
\end{flushright}

\newpage

% ============================================================
% LAB CONTENT
% ============================================================

\section{Learning Outcomes}

After completing this lab students will be able to:

\begin{itemize}
\item Implement the K-Means clustering algorithm from scratch.
\item Understand how different distance metrics affect clustering.
\item Compare clustering results on separable vs overlapping datasets.
\item Visualize cluster decision boundaries.
\end{itemize}

% ------------------------------------------------

\section{Datasets}

In this lab we will experiment with two datasets.

\subsection{Dataset 1: Iris Dataset}

The Iris dataset contains 150 samples.

For visualization purposes we only use two features:

\begin{itemize}
\item Sepal Length
\item Petal Length
\end{itemize}

\subsection{Dataset 2: Synthetic Dataset}

We will also generate a synthetic dataset using:

\begin{verbatim}
sklearn.datasets.make_blobs
\end{verbatim}

Two versions will be created:

\begin{itemize}
\item Well-separated clusters
\item Overlapping clusters
\end{itemize}

This allows us to observe how K-Means behaves when clusters are not clearly separable.

% ------------------------------------------------

\section{Distance Metrics}

In this lab you will experiment with three distance metrics.

\subsection{Euclidean Distance}

\[
d(x,y) = \sqrt{\sum (x_i - y_i)^2}
\]

Produces circular cluster boundaries.

\subsection{Manhattan Distance}

\[
d(x,y) = \sum |x_i - y_i|
\]

Produces diamond-shaped cluster boundaries.

\subsection{Cosine Distance}

Cosine similarity measures the angle between two vectors.

\[
\text{similarity} =
\frac{x \cdot y}{||x||\,||y||}
\]

Cosine distance:

\[
d(x,y) = 1 - similarity
\]

This metric is commonly used in:

\begin{itemize}
\item Text clustering
\item Document similarity
\item Embedding spaces
\end{itemize}

% ------------------------------------------------

\section{Part A: Distance Function (3 Marks)}

\begin{tcolorbox}

Implement the function:

\begin{verbatim}
_distance(x1, x2)
\end{verbatim}

The function must support the following metrics:

\begin{itemize}
\item Euclidean
\item Manhattan
\item Cosine
\end{itemize}

Use NumPy operations only.

\end{tcolorbox}

% ------------------------------------------------

\section{Part B: Cluster Assignment (4 Marks)}

\begin{tcolorbox}

Implement:

\begin{verbatim}
assign_clusters(X, centroids)
\end{verbatim}

Steps:

\begin{itemize}
\item Compute the distance from each point to every centroid
\item Assign the point to the closest centroid
\item Return cluster labels
\end{itemize}

\end{tcolorbox}

% ------------------------------------------------

\section{Part C: Update Centroids (4 Marks)}

\begin{tcolorbox}

Implement:

\begin{verbatim}
update_centroids(X, labels)
\end{verbatim}

Steps:

\begin{itemize}
\item For each cluster compute the mean of all assigned points
\item Return updated centroid coordinates
\end{itemize}

\end{tcolorbox}

% ------------------------------------------------

\section{Part D: Full K-Means Algorithm (4 Marks)}

\begin{tcolorbox}

Implement:

\begin{verbatim}
fit(X)
predict(X)
\end{verbatim}

Algorithm:

\begin{enumerate}
\item Initialize K centroids randomly
\item Assign points to nearest centroid
\item Update centroid locations
\item Repeat until centroids converge or max iterations reached
\end{enumerate}

\end{tcolorbox}

% ------------------------------------------------

\section{Part E: Boundary Visualization (Bonus)}

\begin{tcolorbox}

Create a function:

\begin{verbatim}
plot_decision_boundaries(model, X)
\end{verbatim}

Steps:

\begin{itemize}
\item Create a mesh grid over the feature space
\item Predict cluster label for each grid point
\item Plot colored decision regions
\end{itemize}

Run the visualization for:

\begin{itemize}
\item Euclidean distance
\item Manhattan distance
\item Cosine distance
\end{itemize}

\end{tcolorbox}

% ------------------------------------------------

\section{Part F: Experimental Comparison}

Run K-Means on both datasets.

\begin{center}
\begin{tabular}{|c|c|c|}
\hline
Dataset & Distance Metric & Observations \\
\hline
Iris & Euclidean & \\
\hline
Iris & Manhattan & \\
\hline
Iris & Cosine & \\
\hline
Separable Blobs & Euclidean & \\
\hline
Overlapping Blobs & Euclidean & \\
\hline
\end{tabular}
\end{center}

Discuss:

\begin{itemize}
\item How do distance metrics change cluster shapes?
\item Which metric performs best for overlapping clusters?
\item Does K-Means always converge to the same solution?
\end{itemize}

% ------------------------------------------------

\section{Cluster Quality: Within Cluster Sum of Squares}

K-Means attempts to minimize the following objective function:

\[
J = \sum_{i=1}^{n} ||x_i - \mu_{c_i}||^2
\]

Where:

\begin{itemize}
\item $x_i$ is a data point
\item $\mu_{c_i}$ is the centroid of the cluster assigned to $x_i$
\end{itemize}

This value is often called:

\begin{itemize}
\item Within Cluster Sum of Squares (WCSS) or Sum of Squared Errors (SSE) 
\item Inertia
\end{itemize}

Lower values indicate tighter clusters.

% ------------------------------------------------

\section{Choosing the Number of Clusters: Elbow Method}

The number of clusters $K$ is usually unknown in real applications.

The \textbf{Elbow Method} helps estimate a reasonable value for $K$.

Procedure:

\begin{enumerate}
\item Run K-Means for different values of $K$
\item Compute the inertia for each $K$
\item Plot inertia vs $K$
\end{enumerate}

Typically the curve looks like:

\begin{center}
High decrease $\rightarrow$ elbow $\rightarrow$ slow decrease
\end{center}

The ``elbow'' point suggests the optimal number of clusters.

% ------------------------------------------------

\section{Part F: Compute Inertia (3 Marks)}

\begin{tcolorbox}

Implement the function:

\begin{verbatim}
compute_inertia(X, labels, centroids)
\end{verbatim}

Steps:

\begin{itemize}
\item For every point compute the squared distance to its centroid
\item Sum the squared distances
\item Return the total inertia
\end{itemize}

\end{tcolorbox}

% ------------------------------------------------

\section{Part G: Elbow Method (Bonus)}

\begin{tcolorbox}

Create a function:

\begin{verbatim}
plot_elbow_curve(X, k_values)
\end{verbatim}

Steps:

\begin{itemize}
\item Train K-Means for multiple values of $K$
\item Compute inertia for each value
\item Plot inertia vs $K$
\end{itemize}

Test with:

\begin{verbatim}
k_values = range(1,10)
\end{verbatim}

Which value of $K$ appears optimal for the dataset?

\end{tcolorbox}

\end{document}